Fix \(\mathbf M\in\mathcal C(\mathcal P)\). We verify each transition against the node denotation in the DCTFTR grammar.

For \(\mathsf{NORM}\), every constructor in the checked trace is either consistency, unnesting, contradiction elimination, Boolean normalization, or disjoint finite additivity. Each constructor is a pointwise identity on the finite recursive SCM, so the raw event and its normal form have equal probability. For \(\mathsf{ANCESTOR}\), summing the finitely many omitted assignments gives
\[
P_{M^\star}(e)=\sum_xP_{M^\star}(e,x),\tag{1}
\]
which is exactly the successor-node denotation.

For \(\mathsf{SPLIT}\), condition on the finitely many complete response profiles used by the checked derivation. Recursive compatibility makes the event indicator deterministic and district factorization gives
\[
P_{M^\star}(e)
=\sum_{\omega}\mathbf 1_e(\omega)
\prod_{D'\in\mathcal D(G)}q_{D'}^\star(\omega_{D'}).\tag{2}
\]
The split certificate groups the finite sums and products in (2) into the generated acyclic successor DAG. Its checker uses only finite distributivity, marginalization, and factorization of the original independent district blocks; it does not assert independence merely because an induced subgraph has split. Hence the old and new node denotations agree.

At \(\mathsf{EXTRACT}\), every accepted raw row is certified district-local, and every accepted derived row satisfies the district-cell equality. The admissibility proof is either \(s=\star\) or \(D\cap\Delta_s=\varnothing\), so district-kernel invariance permits the target-law substitution only in those cases. Stacking precisely the accepted equalities, together with normalization, yields
\[
A_Dq_D^\star=m_D.\tag{3}
\]
Rejected candidates contribute no equation. Thus replacing \(\operatorname{NeedRows}(\tau,D,b)\) by the corresponding face node preserves its denotation \(bq_D^\star\).

At \(\mathsf{FACE\_ID}\), the checked decomposition has
\[
b=\lambda^\top A_D+\sum_{\omega\in Z_D(m_D)}\zeta_\omega e_\omega.
\]
Using (3) and the certified forced zeros gives
\[
bq_D^\star=\lambda^\top m_D,\tag{4}
\]
which is the resolved successor expression. A failed membership test terminates without replacing the root value. A ready \(\mathsf{MARGINALIZE}\) step preserves denotation by finite additivity, and a ready \(\mathsf{PRODUCT}\) step preserves it by multiplication of the already certified district factors in (2). Finally, \(\mathsf{CONDITION}\) retains the same numerator and denominator root values and constructs a quotient only when its separate certificate proves strict positivity throughout the compatibility fiber.

For the operational claims, a transition of phase \(j\) removes one node counted in the \(j\)-th rank coordinate, creates no node in an earlier phase, and leaves every earlier coordinate fixed. Later coordinates may increase, so the lexicographic rank nevertheless strictly decreases. The fixed rule priority and canonical least-node order select at most one successor. Conversely, in a well-typed nonterminal configuration, the least unresolved node is in one of the eight declared phases. The corresponding checker either constructs its successor or returns the specified typed failure; all partial-operation failures are explicit branches. Therefore some rule always applies and \(\operatorname{STUCK}\) is unreachable.